\documentclass[11pt,a4paper]{article}

\usepackage[margin=1in]{geometry}
\usepackage[T1]{fontenc}
\usepackage[utf8]{inputenc}
\usepackage{lmodern}
\usepackage{xcolor}
\usepackage{booktabs}
\usepackage{longtable}
\usepackage{array}
\usepackage{enumitem}
\usepackage{amsmath}
\usepackage{listings}
\usepackage[hidelinks,breaklinks=true]{hyperref}

\definecolor{codebg}{RGB}{248,248,248}
\definecolor{codeframe}{RGB}{220,220,220}
\definecolor{keywordblue}{RGB}{0,70,150}
\definecolor{commentgreen}{RGB}{0,110,55}
\definecolor{stringred}{RGB}{150,35,35}

\lstdefinestyle{cppstyle}{
  language=C++,
  basicstyle=\ttfamily\scriptsize,
  keywordstyle=\color{keywordblue}\bfseries,
  commentstyle=\color{commentgreen},
  stringstyle=\color{stringred},
  numbers=left,
  numberstyle=\tiny\color{gray},
  numbersep=6pt,
  breaklines=true,
  keepspaces=true,
  showstringspaces=false,
  backgroundcolor=\color{codebg},
  frame=single,
  rulecolor=\color{codeframe},
  columns=fullflexible,
  captionpos=b
}

\lstdefinestyle{cmdstyle}{
  basicstyle=\ttfamily\scriptsize,
  breaklines=true,
  keepspaces=true,
  showstringspaces=false,
  backgroundcolor=\color{codebg},
  frame=single,
  rulecolor=\color{codeframe},
  columns=fullflexible,
  captionpos=b
}

\newcommand{\codepath}[1]{\path{#1}}
\newcommand{\ms}{\,\mathrm{ms}}

\title{Implementation Report: PBC-Secured Federated Learning in VANET}
\author{VANET ns-3 Security Research Workspace}
\date{May 2, 2026}

\begin{document}
\maketitle

\begin{abstract}
This report documents the implemented federated learning (FL) extension for the
existing VANET security simulation. The implementation integrates hierarchical
FL with the current pairing-based cryptography (PBC) VANET security model. The
implemented flow uses vehicles as local learners, RSUs as edge aggregators, and
the BS as the global aggregator. FL messages are protected using PBC signatures
in the LTE experiment reported here. The report explains the protocol formulas,
implementation mapping, run configuration, generated CSV files, and measured
latency and power results from the LTE/PBC/FL run.
\end{abstract}

\tableofcontents
\newpage

\section{Purpose of the Implementation}

The goal was to extend the existing VANET security simulator so that it can
simulate a secure hierarchical federated learning workflow. The original
security project already implemented vehicle registration, TA certification,
KGC partial-key generation, PBC pseudonym material, PBC signing, PBC aggregate
verification, and analytical power/latency measurement. The new FL extension
adds the following learning workflow on top of that security layer:

\begin{itemize}[leftmargin=*]
  \item Vehicles receive a global model vector from their serving RSU.
  \item Each selected vehicle computes a local model update.
  \item Each selected vehicle applies deterministic pairwise masking.
  \item Each selected vehicle signs its FL update using PBC.
  \item RSUs verify received FL updates and perform edge aggregation.
  \item The BS receives RSU aggregates and performs global aggregation.
  \item FL round metrics, security latency, communication latency, and energy are written to CSV files.
\end{itemize}

The current ns-3 implementation is a lightweight model-vector simulation. It is
not embedding TensorFlow inside ns-3. The real ML baseline remains in
\codepath{Federated-Learning-VANET}, while ns-3 measures networking, security,
latency, and energy behavior.

\section{Implemented Architecture}

The implemented architecture follows a three-layer hierarchy:

\begin{center}
\begin{tabular}{p{0.25\linewidth}p{0.60\linewidth}}
\toprule
\textbf{Layer} & \textbf{Role in FL workflow} \\
\midrule
Vehicles & Local learners. They compute local updates, add masks, and upload signed updates. \\
RSUs & Edge aggregators. They distribute the global model, verify vehicle updates, and aggregate local updates. \\
BS & Global aggregator. It collects RSU aggregates and computes the next global model. \\
\bottomrule
\end{tabular}
\end{center}

The implemented communication path is:

\[
\text{BS} \rightarrow \text{RSU} \rightarrow \text{Vehicle}
\rightarrow \text{RSU} \rightarrow \text{BS}
\]

The implemented security path is:

\[
\text{FL update bytes} \rightarrow \text{PBC signature} \rightarrow
\text{RSU-side verification/aggregation}
\]

\section{Federated Learning Formulas}

For each selected vehicle \(i\), the local dataset is \(D_i\), the current global
model is \(w^{(t)}\), and the local loss is:

\[
F_i(w)=\frac{1}{|D_i|}\sum_{x\in D_i}\ell(w;x)
\]

The local update is represented as:

\[
w_i^{(t+1)} = w^{(t)} - \eta \nabla F_i(w^{(t)})
\]

For secure aggregation, each vehicle applies a mask \(r_i^{(t)}\):

\[
\widetilde{w}_i^{(t+1)} = w_i^{(t+1)} + r_i^{(t)}
\]

The implementation uses deterministic pairwise masks inside each RSU group so
that the weighted masks cancel during edge aggregation:

\[
\sum_{i\in V_r}\frac{|D_i|}{D_r}r_i^{(t)} = 0
\]

where:

\[
D_r = \sum_{j\in V_r}|D_j|
\]

The RSU computes the edge aggregate:

\[
\widetilde{w}_r^{(t+1)}=
\sum_{i\in V_r}\frac{|D_i|}{D_r}\widetilde{w}_i^{(t+1)}
\]

Because the masks cancel, this becomes:

\[
\widetilde{w}_r^{(t+1)}=
\sum_{i\in V_r}\frac{|D_i|}{D_r}w_i^{(t+1)}
\]

The BS computes the global model:

\[
w^{(t+1)}=\sum_{r\in R}\frac{D_r}{\sum_{k\in R}D_k}\widetilde{w}_r^{(t+1)}
\]

Expanded over vehicles:

\[
w^{(t+1)}=\sum_{r\in R}\sum_{i\in V_r}
\frac{|D_i|}{\sum_{j\in V'}|D_j|}w_i^{(t+1)}
\]

\section{PBC Security Integration}

The FL extension reuses the PBC security primitives already implemented in the
VANET security module. For each FL update, the vehicle signs the serialized FL
payload. The signed material includes the active pseudonym, public verification
key, random commitment, and signature component.

The PBC signing equation used by the simulator is:

\[
\psi_i = S_i + x_iX + w_iR_i
\]

where:

\begin{itemize}[leftmargin=*]
  \item \(S_i\) is the KGC-issued partial private key.
  \item \(x_i\) is the vehicle private verification scalar.
  \item \(Q_i=x_iP\) is the corresponding public verification key.
  \item \(w_i\) is per-message signing randomness.
  \item \(W_i=w_iP\) is the random commitment.
  \item \(X=H(\theta)\), where \(\theta\) is the FL round state.
  \item \(R_i=H(\theta\parallel M_i\parallel PID_i\parallel Q_i\parallel W_i)\).
\end{itemize}

The aggregate verification equation is:

\[
e(\Psi,P) \stackrel{?}{=}
e(P_{pub},\sum_i\Pi_i)\cdot e(X,\sum_iQ_i)\cdot
\prod_i e(R_i,W_i)
\]

where:

\[
\Psi=\sum_i\psi_i
\]

\section{Implementation Mapping}

\begin{longtable}{p{0.27\linewidth}p{0.62\linewidth}}
\toprule
\textbf{Implementation item} & \textbf{Implemented responsibility} \\
\midrule
\codepath{vanet-message.h} & Adds FL message types: \codepath{FL_ROUND_START}, \codepath{FL_MODEL_DOWNLOAD}, \codepath{FL_UPDATE_UPLOAD}, \codepath{FL_RSU_AGGREGATE}, \codepath{FL_GLOBAL_MODEL}, and \codepath{FL_ROUND_DONE}. \\
\codepath{fl-vanet-apps.h/.cc} & Implements vehicle-side local training, masking, PBC/ECC signing, RSU update verification, RSU edge aggregation, BS global aggregation, and FL CSV logging. \\
\codepath{vanet-security-lte.cc} & Adds LTE command-line flags and installs FL BS, RSU, and vehicle apps. \\
\codepath{vanet-security-nr.cc} & Adds the same FL support for NR/6G-like experiments. \\
\codepath{vanet-power-model.cc} & Reused to record computation latency, communication airtime, security energy, communication energy, and stage summaries. \\
\codepath{Federated-Learning-VANET/train_fl_baseline.py} & Provides the external Python ML baseline for real model accuracy, precision, recall, and loss. \\
\bottomrule
\end{longtable}

\subsection{Core Implementation Sketch}

\begin{lstlisting}[style=cppstyle,caption={Simplified implementation sketch of the FL/PBC flow.}]
// BS round start
payload = Serialize(round, globalModel);
SendToEachRsu(FL_MODEL_DOWNLOAD, payload);

// Vehicle local update
localModel = globalModel;
localModel[d] -= learningRate * SyntheticGradient(vehicleId, round, d);
maskedUpdate = localModel + PairwiseMask(vehicleId, rsuGroup, round);

// Vehicle PBC signing
pid = EncodePid(pid1, pid2, tiUs);
theta = "fl-round-" + round;
SignMessage(partialKey, secretXi, pid, theta, payload, wi, psi);
SendToRsu(FL_UPDATE_UPLOAD, payload, pbcAuthHeader);

// RSU aggregation
VerifyAggregate(pidList, qiList, wiList, messageList, theta, aggregatePsi);
edgeModel = Sum_i((datasetSize_i / rsuDatasetTotal) * update_i);
SendToBs(FL_RSU_AGGREGATE, edgeModel);

// BS global aggregation
globalModel = Sum_r((rsuDatasetTotal_r / totalDataset) * edgeModel_r);
WriteFlCsvRow(round);
\end{lstlisting}

\section{Experiment Configuration}

The reported experiment used the following command:

\begin{lstlisting}[style=cmdstyle,caption={LTE PBC federated learning run command.}]
LSAN_OPTIONS=detect_leaks=0 ./ns3 run "vanet-security-lte \
  --securityMode=pbc \
  --enableFl=true \
  --flSecurityMode=pbc \
  --flRounds=3 \
  --flParticipants=10 \
  --flModelDim=16 \
  --simTime=30 \
  --progressLog=true \
  --outCsv=results/fl/lte_pbc_fl_power.csv \
  --stageCsv=results/fl/lte_pbc_fl_stages.csv \
  --flCsv=results/fl/lte_pbc_fl.csv \
  --flStageCsv=results/fl/lte_pbc_fl_stage_summary.csv"
\end{lstlisting}

The run configuration was:

\begin{center}
\begin{tabular}{p{0.35\linewidth}p{0.45\linewidth}}
\toprule
\textbf{Parameter} & \textbf{Value} \\
\midrule
Network label & \codepath{5g_lte} \\
Security mode & \codepath{pbc} \\
Simulation time & 30 seconds \\
FL rounds requested & 3 \\
FL participants & 10 vehicles \\
RSU count & 2 \\
Model dimension & 16 parameters \\
FL masking & Enabled \\
Power model & Analytical CPU/radio energy model \\
\bottomrule
\end{tabular}
\end{center}

\section{Generated Output Files}

\begin{longtable}{p{0.35\linewidth}p{0.55\linewidth}}
\toprule
\textbf{File} & \textbf{Purpose} \\
\midrule
\codepath{results/fl/lte_pbc_fl.csv} & Completed FL round summary. A row is written only when a BS global aggregation round completes. \\
\codepath{results/fl/lte_pbc_fl_power.csv} & Periodic VANET-wide power, latency, infection, registration, KGC, and TA metrics. \\
\codepath{results/fl/lte_pbc_fl_stages.csv} & Stage-level latency and energy for security, communication, and FL stages. \\
\codepath{results/fl/lte_pbc_fl_stage_summary.csv} & Same stage-level summary written under the FL security label. \\
\codepath{results/fl/lte_pbc_fl_report.md} & Human-readable Markdown report generated from the CSV files. \\
\bottomrule
\end{longtable}

\section{Completed FL Round Results}

The main FL result row is shown in Table~\ref{tab:fl-round}. The row proves that
one complete global FL round finished successfully.

\begin{table}[h]
\centering
\caption{Completed LTE/PBC/FL round.}
\label{tab:fl-round}
\begin{tabular}{lr}
\toprule
\textbf{Metric} & \textbf{Value} \\
\midrule
Completion time & 10.402010 s \\
Round completed & 0 \\
Network label & 5g\_lte \\
Security mode & pbc \\
Selected vehicles & 10 \\
RSUs & 2 \\
Updates sent & 10 \\
Updates verified & 10 \\
Updates rejected & 0 \\
RSU edge aggregates sent & 2 \\
Global round completed & 1 \\
Model download bytes & 1632 \\
Update upload bytes & 1440 \\
Loss proxy & 0.500000 \\
Accuracy proxy & 0.500000 \\
\bottomrule
\end{tabular}
\end{table}

The most important correctness result is:

\[
\text{updates\_sent}=10,\quad
\text{updates\_verified}=10,\quad
\text{updates\_rejected}=0
\]

This means every submitted FL update passed the PBC-secured verification path in
this clean run.

\section{FL Latency and Energy Results}

\begin{table}[h]
\centering
\caption{FL round latency and energy at global completion time.}
\label{tab:fl-latency-energy}
\begin{tabular}{lr}
\toprule
\textbf{Metric} & \textbf{Value} \\
\midrule
Average local training latency & 0.002526 ms \\
Average update signing latency & 18.265342 ms \\
Average update verification latency & 89.195864 ms \\
Average RSU edge aggregation latency & 0.002608 ms \\
Average BS global aggregation latency & 0.007223 ms \\
Security energy at FL completion & 92.456386 J \\
Communication energy at FL completion & 2.002021 J \\
Total energy at FL completion & 94.458407 J \\
Average power at FL completion & 9.080784 W \\
\bottomrule
\end{tabular}
\end{table}

The latency results show that local training, masking, edge aggregation, and
global aggregation are very small in the lightweight ns-3 FL model. The dominant
operations are PBC signing and PBC aggregate verification.

\section{Final VANET-Wide Power and Delay Results}

The final periodic CSV row was recorded at simulated time 29 seconds. These
values describe the entire LTE/PBC/FL simulation state near the end of the run.

\begin{table}[h]
\centering
\caption{Final VANET-wide LTE/PBC/FL metrics at 29 seconds.}
\label{tab:final-power}
\begin{tabular}{lr}
\toprule
\textbf{Metric} & \textbf{Value} \\
\midrule
Infected ratio & 1.000000 \\
Propagation distance & 103.2 \\
Verification failures & 0 \\
Average V2V delay & 1113.69 ms \\
Average V2I uplink delay & 18.927 ms \\
Average V2I downlink delay & 1368.4 ms \\
Average V2I RTT & 1382.54 ms \\
Average registration delay & 495.8 ms \\
Average KGC RTT & 8.07766 ms \\
Average TA RTT & 8.06223 ms \\
Security energy & 312.084 J \\
Communication energy & 5.54503 J \\
Total energy & 317.629 J \\
Average security power & 10.7615 W \\
Average communication power & 0.191208 W \\
Average total power & 10.9527 W \\
Average signing latency & 22.2315 ms \\
Average aggregation latency & 0.173781 ms \\
Average verification latency & 98.1813 ms \\
Average partial-key latency & 5.31927 ms \\
\bottomrule
\end{tabular}
\end{table}

The final result shows that security computation dominates total energy:

\[
E_{security}=312.084\,J,\quad E_{communication}=5.54503\,J
\]

Thus:

\[
\frac{E_{security}}{E_{total}}\approx
\frac{312.084}{317.629}=98.25\%
\]

This confirms that, in this PBC-secured LTE experiment, cryptographic processing
is much more expensive than analytical radio communication energy.

\section{FL Stage-Level Results}

\begin{longtable}{p{0.25\linewidth}p{0.20\linewidth}r r r}
\caption{FL-specific communication and computation stages.}\label{tab:fl-stages}\\
\toprule
\textbf{Stage} & \textbf{Role} & \textbf{Count} & \textbf{Avg. ms} & \textbf{Energy J} \\
\midrule
\endfirsthead
\toprule
\textbf{Stage} & \textbf{Role} & \textbf{Count} & \textbf{Avg. ms} & \textbf{Energy J} \\
\midrule
\endhead
communication\_rx & bs:fl\_rsu\_aggregate\_upload & 3 & 0.014160 & 0.000212 \\
communication\_rx & rsu:fl\_model\_download & 4 & 0.012880 & 0.000258 \\
communication\_rx & rsu:fl\_update\_upload & 19 & 0.057200 & 0.005434 \\
communication\_rx & vehicle:fl\_model\_download & 20 & 0.012880 & 0.000206 \\
communication\_tx & bs:fl\_model\_download & 4 & 0.012880 & 0.000052 \\
communication\_tx & rsu:fl\_model\_download & 20 & 0.012880 & 0.000258 \\
communication\_tx & rsu:fl\_rsu\_aggregate\_upload & 3 & 0.014160 & 0.000042 \\
communication\_tx & vehicle:fl\_update\_upload & 19 & 0.057200 & 0.000217 \\
fl\_edge\_aggregate & rsu & 3 & 0.002431 & 0.000109 \\
fl\_global\_aggregate & bs & 1 & 0.007223 & 0.000108 \\
fl\_local\_train & vehicle & 19 & 0.002572 & 0.000122 \\
fl\_mask\_generate & vehicle & 19 & 0.001457 & 0.000069 \\
\bottomrule
\end{longtable}

The stage counts show that round 0 completed and round 1 partially progressed.
A completed round with 10 vehicles and 2 RSUs requires 10 vehicle updates, 2 RSU
edge aggregates, and 1 BS global aggregate. The CSV contains:

\[
\text{fl\_global\_aggregate count}=1
\]

Therefore, exactly one complete BS global FL round was recorded. The larger
counts for model downloads and update uploads indicate that the next FL round
started before the simulation ended.

\section{PBC Security Stage Results}

\begin{longtable}{p{0.30\linewidth}p{0.17\linewidth}r r r}
\caption{PBC security stages during LTE/PBC/FL run.}\label{tab:pbc-stages}\\
\toprule
\textbf{Stage} & \textbf{Role} & \textbf{Count} & \textbf{Avg. ms} & \textbf{Energy J} \\
\midrule
\endfirsthead
\toprule
\textbf{Stage} & \textbf{Role} & \textbf{Count} & \textbf{Avg. ms} & \textbf{Energy J} \\
\midrule
\endhead
pbc\_aggregate & bs\_relay & 215 & 0.172608 & 0.556660 \\
pbc\_full\_key & vehicle & 60 & 1.792057 & 0.268808 \\
pbc\_partial\_key & kgc & 60 & 5.319274 & 4.787346 \\
pbc\_pid2 & ta & 60 & 1.691061 & 1.521955 \\
pbc\_sign & vehicle & 49 & 22.315887 & 2.733696 \\
pbc\_verify\_aggregate & bs\_relay & 215 & 97.382770 & 314.059434 \\
\bottomrule
\end{longtable}

The most expensive security operation is PBC aggregate verification:

\[
\text{avg latency}=97.382770\ms,
\quad \text{energy}=314.059434\,J
\]

This is the main reason that security energy dominates communication energy in
this run. The PBC stage labels are shared by the existing VANET PBC relay path
and the FL update verification path, so Table~\ref{tab:pbc-stages} should be
read as the total PBC workload during the LTE/PBC/FL simulation.

\section{Interpretation of the Results}

The run demonstrates four important outcomes.

\subsection{The FL Pipeline Completed Correctly}

The completed FL row shows that 10 updates were sent, 10 were verified, and 0
were rejected. Both RSUs sent their edge aggregates, and the BS completed one
global aggregation round. This proves that the vehicle--RSU--BS hierarchical FL
pipeline is active and connected to the PBC security layer.

\subsection{Security Verification Is the Main Cost}

The average update signing latency at FL completion was 18.265342 ms, while the
average update verification latency was 89.195864 ms. The final stage summary
also shows that PBC aggregate verification consumed 314.059434 J, which is much
larger than the FL aggregation computation itself. This means PBC security is
computationally expensive, especially at the verification side.

\subsection{Communication Energy Is Small Compared with Security Energy}

At the end of the run, security energy was 312.084 J and communication energy
was 5.54503 J. Thus, almost all measured energy came from cryptographic
processing. This is an important research finding because it shows that, for
PBC-secured VANET-FL, optimizing verification cost may be more important than
optimizing packet airtime.

\subsection{The FL Model Is Lightweight in ns-3}

The local training and aggregation latencies are very small because the ns-3 FL
implementation intentionally uses lightweight numeric vectors. This allows the
simulation to focus on communication, PBC security, latency, and power. Real ML
accuracy, precision, recall, and loss are handled by the Python baseline in
\codepath{Federated-Learning-VANET}.

\section{Limitations and Notes}

\begin{itemize}[leftmargin=*]
  \item The ns-3 FL implementation uses synthetic gradients and model vectors; it does not train TensorFlow models inside ns-3.
  \item The reported accuracy and loss in ns-3 are proxy values used to indicate FL progress, not real ML accuracy.
  \item The PBC stage labels include both VANET PBC relay traffic and FL PBC update traffic.
  \item Only one full FL round completed in this 30-second run. Later rounds partially progressed but did not complete before simulation stop.
  \item To complete all requested three FL rounds, the simulation should be run for a longer duration such as 90 seconds.
\end{itemize}

\section{Recommended Next Run}

For a cleaner multi-round professor demo, run:

\begin{lstlisting}[style=cmdstyle,caption={Recommended three-round LTE/PBC/FL run.}]
LSAN_OPTIONS=detect_leaks=0 ./ns3 run "vanet-security-lte \
  --securityMode=pbc \
  --enableFl=true \
  --flSecurityMode=pbc \
  --flRounds=3 \
  --flParticipants=10 \
  --flModelDim=16 \
  --simTime=90 \
  --progressLog=true \
  --outCsv=results/fl/lte_pbc_fl_power_90s.csv \
  --stageCsv=results/fl/lte_pbc_fl_stages_90s.csv \
  --flCsv=results/fl/lte_pbc_fl_90s.csv \
  --flStageCsv=results/fl/lte_pbc_fl_stage_summary_90s.csv"
\end{lstlisting}

Then verify completed FL rounds with:

\begin{lstlisting}[style=cmdstyle,caption={Commands to check FL completion.}]
cat results/fl/lte_pbc_fl_90s.csv
grep -E "fl_edge_aggregate|fl_global_aggregate|fl_rsu_aggregate_upload" \
  results/fl/lte_pbc_fl_stages_90s.csv
\end{lstlisting}

\section{Conclusion}

The LTE/PBC/FL implementation successfully integrates federated learning with
the existing VANET PBC security model. The completed run proves that vehicles
can receive a global model, compute masked local updates, sign the updates,
upload them to RSUs, and contribute to a BS-level global aggregation round. The
results also show that PBC verification dominates both latency and energy. This
makes PBC aggregate verification the main optimization target for future secure
VANET federated learning experiments.

\end{document}
